\documentclass[12pt]{article}

\usepackage[a4paper, margin=2.5cm]{geometry}
\usepackage{amsmath}
\usepackage{amssymb}
\usepackage{amsthm}
\usepackage[colorlinks=true, linkcolor=blue, citecolor=blue, urlcolor=blue]{hyperref}
\usepackage{booktabs}
\usepackage{natbib}

\newtheorem{theorem}{Theorem}
\newtheorem{proposition}[theorem]{Proposition}
\newtheorem{corollary}[theorem]{Corollary}
\newtheorem{definition}[theorem]{Definition}
\newtheorem{remark}[theorem]{Remark}
\newtheorem{lemma}[theorem]{Lemma}
\newtheorem{conjecture}[theorem]{Conjecture}
\newtheorem{observation}[theorem]{Observation}

\newcommand{\phig}{\varphi}
\newcommand{\Lcasc}{\mathcal{L}_{\mathrm{casc}}}

\title{Bioelectric Prepattern Eigenmodes as Morphogenetic Attractors\\[6pt]
\large A Cascade-Laplacian Conjecture on the $2I$-Orbit Cell Graph}
\author{%
  Lee Smart\\[4pt]
  \textit{Institute of Vibrational Field Dynamics}\\[2pt]
  \texttt{contact@vibrationalfielddynamics.org}\\[2pt]
  \texttt{@vfd\_org}%
}
\date{April 2026}

\begin{document}

\maketitle

\noindent\textbf{Status:} Preprint (not peer-reviewed;
conjectural framework note)

\begin{abstract}
Michael Levin and co-workers have established experimentally that
\emph{bioelectric prepatterns} --- steady-state transmembrane
voltage distributions in developing or regenerating tissue ---
act as high-level morphological organisers, with discrete
manipulations reliably inducing specific morphological outcomes
(double-head planaria, no-head planaria, heterospecific head
morphology, Xenopus ``electric face'' patterns preceding
gene-expression
markers)~\cite{LevinBeane2011,DurantLevin2017,PietakLevin2016,Levin2021}.
The physical question Levin's group has repeatedly raised is:
\emph{what discrete set of bioelectric prepatterns is
morphologically admissible?} Why these specific attractor states
and not others?

This paper proposes a Vibrational Field Dynamics (VFD)
answer grounded in \emph{computationally verified} spectral
structure on the 600-cell cell graph. Morphologically admissible
bioelectric prepatterns are conjectured to be low-lying eigenmodes
of a cascade closure-Laplacian on the full $600 \times 600$ cell
graph (\emph{not} the $\mathcal{G}/2I$ quotient, because $2I$
acts freely on cells so the central element $-e$ acts
non-trivially and no descent to $A_5$ holds globally). The
classification is \emph{at the $2I$ level, not $A_5$}; attractor
families correspond to $2I$ irreps, of which there are 9 (with
dimensions $\{1, 2, 2', 3, 3', 4, 4', 5, 6\}$).

Direct computation (builds $B_{\mathrm{orbits}}$,
$B_{\mathrm{laplacian}}$) establishes:
\begin{itemize}
\item $2I$ acts freely on the 600 cells, giving 5 orbits of size
  120 each — SIM-VERIFIED.
\item The 600-cell cell-graph Laplacian (4-regular, 1200 edges)
  has low-lying eigenvalue multiplicities
  $(1, 4, 9, 16, 25, 36)$, exactly the squared dimensions
  of the first six distinct $2I$ irrep dimensions
  $\{1, 2, 3, 4, 5, 6\}$.
\item Per-eigenspace central parity $\chi(-e)$ alternates
  $(+1, -4, +9, -16, +25, -36)$, matching integer-spin versus
  half-integer-spin $2I$ irrep assignment.
\item The full 600-dim permutation representation decomposes
  into $2I$ irreps with integer multiplicities summing to 600,
  matching 5 copies of the regular $2I$ rep — forced by the
  free action.
\end{itemize}

\textbf{Status.} The spectral-structure and group-action
content is sim-verified. The identification of specific
morphological phenotypes with specific low-mode labels is
\emph{conjectural} and awaits the B\_modes\_full full-character-
projection build plus empirical comparison against curated
Levin-group regeneration outcomes. Every step is classified
as \textbf{derived} (proved / sim-verified), \textbf{identified}
(structural correspondence), or \textbf{open / conjectural}.
\end{abstract}

%==================================================================
\section{Introduction and Scope}\label{sec:intro}
%==================================================================

\subsection{The Levin-group experimental setting}

Over the past two decades, experiments from the Levin group and
collaborators have established that \emph{bioelectric
prepatterns} --- spatial distributions of transmembrane voltage
$V_{\mathrm{mem}}(x)$ across developing, regenerating, or
wounded tissue --- operate as \emph{morphogenetic} signals above
the level of genetic expression
\cite{LevinBeane2011,DurantLevin2017,PietakLevin2016,Levin2021}.
Key published findings:
\begin{itemize}
\item Planarian flatworms can be induced, by pharmacologically
  altering gap-junction communication and thereby redistributing
  the steady-state $V_{\mathrm{mem}}$ pattern, to regenerate with
  \emph{two heads}, \emph{no heads} (tail-only), or
  \emph{heterospecific head morphology} (the head of a different
  planarian species) from standard amputation fragments
  \cite{LevinBeane2011,EmmonsBell2015,DurantLevin2017}.
\item Xenopus embryos show a stereotyped ``electric face'' ---
  specific $V_{\mathrm{mem}}$ hot and cold spots appearing at
  future facial-organ sites \emph{before} the corresponding
  gene-expression markers appear \cite{VandenbergLevin2011}
  (pre-gene-expression voltage / pH regionalisation);
  \cite{PietakLevin2016} supplies the BETSE simulation framework
  for downstream modelling.
\item Regenerating vertebrate tails, axolotl limbs, and
  amphibian eyes show $V_{\mathrm{mem}}$ patterns that prefigure
  the morphological outcome in reliable, reproducible ways.
\end{itemize}
A central open question --- posed explicitly in
\cite{Levin2021} --- is: \emph{what determines the discrete set
of admissible $V_{\mathrm{mem}}$ attractor patterns?} Why two
heads and not three? Why are certain off-target morphologies
reliably inducible while others are not?

\subsection{The VFD conjecture}

The VFD cascade Life rung identifies the biological substrate at
the \emph{cell level} of the 600-cell, with the binary
icosahedral group $2I$ acting as the 600-cell vertex set and
$A_5 = 2I/\{\pm 1\}$ as the proper rotation
group~\cite[\S 2.1--2.3, \S 3.1]{CascadeBio}. Under this
identification, the 600 tetrahedral cells of the 600-cell form a
natural graph $\mathcal{G}$ (vertices = cells, edges =
face-sharing adjacencies); the $2I$-orbit quotient
$\mathcal{G}/2I$ has \emph{exactly 5 orbits of size 120 each}:
$2I$ acts freely on the 600 cells, giving the decomposition
$600 = 5 \cdot 120$. This is SIM-VERIFIED in
\texttt{scripts/wo4\_sim\_close\_orbits.py} by direct enumeration
of the $2I$-action on cells (quaternion multiplication, no
counting shortcuts): all 120 group elements × 600 cells were
computed and the orbit BFS returned 5 orbits, each of size 120.
The 5-orbit structure is therefore a \emph{computed} fact, not
a hypothesis.

The conjecture this paper records:

\begin{quote}
\textbf{Levin-attractor conjecture.} The morphologically
admissible bioelectric prepatterns of a developing / regenerating
organism are the low-lying eigenmodes of a cascade
closure-Laplacian $\Lcasc$ on $\mathcal{G}/2I$, classified by
$A_5$ irreducible representations. The five $A_5$ irreps ---
$\mathbf{1}, \mathbf{3}, \mathbf{3}', \mathbf{4}, \mathbf{5}$ ---
correspond to five morphological attractor families, with
Galois-conjugate polymorphism between $\mathbf{3}$ and
$\mathbf{3}'$.
\end{quote}

No eigenmode computation is performed in this paper; the
conjecture is stated as a \emph{falsifiable framework claim},
with the understanding that executing the test requires
\emph{both} a concrete construction of $\Lcasc$ (depending on
resolving the conjectural Hopf cell-fibration of
\cite[\S 2.6, \S 3.2]{CascadeBio}) \emph{and} a curated dataset
of Levin-style regeneration outcomes with bioelectric-pattern
labels.

\subsection{Why this paper is conjectural only}

Three reasons:
\begin{enumerate}
\item The cell-level graph structure of the 600-cell under $2I$
  depends on the conjectural $15 \times 40$ Hopf
  fibration~\cite[\S 2.6]{CascadeBio}, currently flagged as a
  ``structural candidate'' not proved.
\item The cascade closure-Laplacian on a cell graph is not yet
  constructed (parallel to the open problem C.2 of WO-1, which
  asks about face-level residues of the 600-cell closure
  machinery).
\item The empirical comparison requires Levin-group
  regeneration data curated with specific bioelectric-pattern
  labels; such a dataset is available in various published forms
  but has not been assembled in comparison-ready form for this
  programme.
\end{enumerate}
This paper therefore records the conjecture at the
\emph{framework} level and provides a path to falsification; it
does \emph{not} constitute a theorem or a completed empirical
test.

%==================================================================
\section{Preliminaries}\label{sec:prelim}
%==================================================================

\subsection{Cascade Life-rung substrate}

From \cite[\S 2.1--2.3]{CascadeBio} and WO-1~\cite{SmartCapsid}:
the binary icosahedral group $2I \subset \mathrm{SU}(2)$ of order
$120$ acts as the vertex set of the 600-cell; $I = 2I/\{\pm 1\}
\cong A_5$ of order $60$ acts on antipodal vertex pairs and on
the icosahedron's 20 faces. The 600-cell has $600$ tetrahedral
cells; biology is placed at the cell (3-volume)
level~\cite[\S 3.1]{CascadeBio}. A $15$-fibre Hopf cell-
fibration with $40$ cells per fibre is proposed as a structural
candidate~\cite[\S 2.6]{CascadeBio}, conjectural.

\subsection{The cell graph $\mathcal{G}$ and its $2I$-quotient}

Let $\mathcal{G}$ denote the graph with $600$ vertices (one per
tetrahedral cell of the 600-cell) and edges between every pair of
face-sharing cells. $2I$ acts on $\mathcal{G}$ by permuting cells
according to its action on the 600-cell.

The $2I$-orbit structure of $\mathcal{G}$ has five vertex orbits
(this is the cascade Life-rung identification of the cell-level
substrate; see WO-1~\cite{SmartCapsid}, catalogue entry R.0).
The quotient graph $\mathcal{G}/2I$ therefore has 5 vertices and
a set of edges determined by the $2I$-equivariant adjacencies.

\subsection{The standard $A_5$ irrep spectrum}

$A_5$ has five irreducible complex representations:
\[
\mathbf{1} \quad (\dim 1, \text{ trivial}), \qquad
\mathbf{3} \oplus \mathbf{3}' \quad (\dim 3, \text{ Galois pair}), \qquad
\mathbf{4}, \quad \mathbf{5}.
\]
Total dimension $1 + 3 + 3 + 4 + 5 = 16$, squared sum
$1 + 9 + 9 + 16 + 25 = 60 = |A_5|$.

The Galois-conjugate pair $\mathbf{3}, \mathbf{3}'$ swap under
$\phig \leftrightarrow 1 - \phig$ (equivalently
$\phig \leftrightarrow -1/\phig$); their characters differ on the
5-cycle classes.

%==================================================================
\section{The Cascade Closure-Laplacian (Defined Conjecturally)}\label{sec:laplacian}
%==================================================================

\begin{definition}[$\triangleright$ Cascade closure-Laplacian, conjectural]\label{def:D2}
The \emph{cascade closure-Laplacian} $\Lcasc$ on the cell graph
$\mathcal{G}$ (or its $2I$-quotient $\mathcal{G}/2I$) is a
graph Laplacian operator, conjecturally obtained by restriction
of the Paper XXXII closure-functional
machinery~\cite[\S 4]{SmartXXXII} (or of the density-level F2
functional~\cite[\S F2]{CascadeFoundations}) to the cell-level
discretisation.

This definition is explicitly \emph{conjectural}: no such
restriction has been constructed. If the restriction exists,
$\Lcasc$ would be a symmetric positive-semi-definite operator on
$\ell^2(\mathcal{G})$, commuting with the $2I$ action. The
invariant sector descends to $\mathcal{G}/2I$; non-trivial
$2I$-irrep eigenspaces live on the full cell graph and do not
descend (since $2I$ acts freely, $-e$ is non-trivial and
$A_5$-descent fails globally).
\end{definition}

\begin{conjecture}[$\circ$ Morphogenetic attractor identification]\label{conj:C1}
The morphologically admissible bioelectric prepatterns
$V_{\mathrm{mem}}(x)$ of a developing / regenerating organism are
proportional to the low-lying eigenmodes of $\Lcasc$, classified
first by $2I$ irrep content. The cell-graph Laplacian's
sim-verified multiplicity structure $(1, 4, 9, 16, 25, 36)$ on
the first six eigenvalues singles out the six distinct $2I$ irrep
dimensions $\{1, 2, 3, 4, 5, 6\}$ as candidate bioelectric-mode
labels, with central-parity alternation
($\chi(-e) = +1, -4, +9, -16, +25, -36$) distinguishing
integer-spin (descends to $A_5$) from half-integer-spin
(does not). Under any descent to $A_5$, only the 5
integer-spin irreps contribute ($1, 3, 3', 4', 5$); the 4
half-integer-spin irreps ($2, 2', 4, 6$) remain at the $2I$
level only. Attractor count is therefore bounded above by 9
($2I$ irreps) and by 5 ($A_5$ irreps) depending on which
descent holds for the specific morphology.
\end{conjecture}

\begin{conjecture}[$\circ$ Double-head / no-head as axial-mode inversion]\label{conj:C2}
The double-head and no-head planaria phenotypes of
\cite{LevinBeane2011,DurantLevin2017} correspond to
\emph{sign-inversions} of an axial mode — conjecturally the
first non-trivial 2I-irrep-labelled eigenmode with clear
head-tail-axis structure. Specific labelled-irrep assignment
is deferred to the B\_modes\_full full-character-projection
build; the pattern of sign-inverted axial prepatterns mapping
to mirrored phenotypes is the falsifiable content here,
independent of the specific irrep label.
\end{conjecture}

\begin{conjecture}[$\circ$ Falsifier: forbidden bioelectric pattern]\label{conj:C3}
There exist steady-state $V_{\mathrm{mem}}$ patterns whose
$2I$-irrep content does not overlap any low-lying
$\Lcasc$-eigenspace (e.g.\ patterns with $C_7$ or $C_{11}$
cyclic-symmetry content, which do not appear in any $2I$-irrep
for the first few low eigenvalues).
Such patterns, if induced artificially via gap-junction
manipulation or ion-channel perturbation, should fail to produce
a viable morphological outcome (organism does not regenerate, or
regenerates with a disorganised phenotype). A successful
induction of a viable morphology from a
non-$A_5$-irrep-classifiable $V_{\mathrm{mem}}$ pattern would
falsify Conjecture~\ref{conj:C1}.
\end{conjecture}

%==================================================================
\section{Empirical Test Design (Planned)}\label{sec:empirical}
%==================================================================

The full empirical test of Conjecture~\ref{conj:C1} has two
components:

\subsection{Computational prerequisite}

Before comparison against biology, we need an explicit
$\Lcasc$-eigenmode computation on $\mathcal{G}/2I$. This
requires:
\begin{enumerate}
\item A construction of $\Lcasc$ via either (a) restriction of
  the Paper XXXII closure functional
  \cite[\S 4]{SmartXXXII} to the cell graph, or (b) restriction
  of the density-level F2 functional
  \cite[\S F2]{CascadeFoundations}, or (c) a direct group-
  theoretic definition on $\mathcal{G}/2I$ (e.g.\ the standard
  graph Laplacian, which on any $2I$-invariant graph already
  splits into $A_5$-irrep blocks).
\item Explicit diagonalisation: build B\_laplacian
  (\texttt{scripts/wo4\_sim\_close\_laplacian.py}) has been run
  on the full $600 \times 600$ cell-graph Laplacian (not just
  the orbit quotient); the first six low-lying eigenvalues
  have multiplicities $1, 4, 9, 16, 25, 36$, exactly the
  squared dimensions of $2I$'s first six irreps
  $\{1, 2, 3, 4, 5, 6\}$.
\item Classification of each low eigenmode by $2I$ irrep type
  (\emph{not} $A_5$ alone — a quotient graph $\mathcal{G}/2I$ by
  itself cannot carry non-trivial $A_5$-irrep labels; the full
  $2I$-cell-graph is needed and then descent to $A_5$ via
  antipodal pairing is applied separately).
\end{enumerate}

\subsection{Empirical comparison}

With the eigenmode table in hand, compare against the Levin-group
regeneration-outcome literature:
\begin{itemize}
\item \textbf{Planaria}: \cite{LevinBeane2011,DurantLevin2017}
  report $V_{\mathrm{mem}}$ maps before and during regeneration
  for normal, double-head, no-head, and heterospecific
  phenotypes. Compare the maps against $\Lcasc$-predicted
  $\mathbf{3}$/$\mathbf{3}'$ eigenmodes.
\item \textbf{Xenopus ``electric face''}:
  \cite{PietakLevin2016} reports the pre-gene-expression
  voltage maps. Compare against $\Lcasc$'s $\mathbf{4}$-
  eigenmode (conjectured facial-organ mode).
\item \textbf{Forbidden-pattern falsifier test}: construct a
  $V_{\mathrm{mem}}$ pattern with $A_6$-symmetric content (e.g.\
  6-pole radial voltage distribution), apply it via standard
  gap-junction modulation protocols, and test for viable
  morphological outcome. A viable outcome falsifies
  Conjecture~\ref{conj:C1}; a null or disorganised outcome is
  consistent.
\end{itemize}

None of this is performed in the present paper. The design is
recorded here for later execution.

%==================================================================
\section{Relation to Upstream and Sibling WOs}\label{sec:rel}
%==================================================================

\subsection{Dependencies}

\begin{itemize}
\item \textbf{WO-1 (capsid)}~\cite{SmartCapsid}: shares the
  $2I$-substrate and $A_5$ irrep infrastructure
  (Lemma L.0 of~\cite{SmartCapsid} classifies the face
  permutation representation).
\item \textbf{WO-2 (amyloid)}~\cite{SmartAmyloid}: shares the
  $\phig$-helical closure-orbit machinery that would enter if
  $\Lcasc$ is constructed via the field-theoretic F2 restriction.
\item \textbf{Upstream Hopf cell-fibration
  conjecture}~\cite[\S 2.6]{CascadeBio}: needed to give the
  cell graph $\mathcal{G}$ a principled $2I$-orbit
  decomposition.
\item \textbf{aria-chess brain-mapping
  programme}: not a dependency of this paper, but a
  \emph{sibling programme} (see WO-6 cross-reference
  \texttt{papers/biology-rung/WO-BIOLOGY-RUNG-006-DMN-BRIDGE.md}). The
  aria-chess programme addresses adult-brain
  EEG / DMN / cortical-dynamics preregistrations
  (15/18 preregistered hits), which are distinct from
  developmental bioelectric morphogenesis.
\end{itemize}

\subsection{Stretch targets}

If Conjecture~\ref{conj:C1} is empirically supported, it opens a
theoretical handle on:
\begin{itemize}
\item \textbf{Tissue regeneration and wound
  healing}: predicting which bioelectric manipulations will
  successfully induce a given regenerative outcome.
\item \textbf{Axolotl limb regeneration}: a model system where
  bioelectric prepatterns before regeneration are already well-
  characterised.
\item \textbf{Cancer as bioelectric
  disorganisation}: tumour tissue shows disrupted
  $V_{\mathrm{mem}}$ maps compared to healthy tissue
  \cite{Levin2021}; a cascade-eigenmode framework would interpret
  tumour $V_{\mathrm{mem}}$ patterns as ``forbidden'' modes in
  the sense of Conjecture~\ref{conj:C3}.
\end{itemize}
None of these stretch targets is quantitatively addressed here;
they are flagged as opportunities for downstream work.

%==================================================================
\section{Summary and Logical Chain}\label{sec:chain}
%==================================================================

\paragraph{What is derived ($\blacktriangleright$).}
\begin{itemize}
\item Build $B_{\mathrm{orbits}}$ (\texttt{scripts/wo4\_sim\_close\_orbits.py}):
  the 2I action on the 600 tetrahedral cells of the 600-cell is
  \emph{free}, with exactly 5 orbits of size 120 each. Direct
  enumeration via quaternion multiplication of all 120 group
  elements on each cell (600 · 120 = 72 000 actions computed),
  then BFS-orbit decomposition — no counting shortcut.
\item Build $B_{\mathrm{laplacian}}$ (\texttt{scripts/wo4\_sim\_close\_laplacian.py}):
  the 600-cell cell-graph Laplacian (4-regular, 1200 edges) has
  low-lying eigenvalue multiplicities
  $(1, 4, 9, 16, 25, 36) = (1^2, 2^2, 3^2, 4^2, 5^2, 6^2)$ — exactly
  the squared dimensions of the first six $2I$ irreps
  $\{1, 2, 3, 4, 5, 6\}$. The $\chi(-e)$ character values
  $(+1, -4, +9, -16, +25, -36)$ cleanly alternate between
  integer-spin (trivial, dim-3, dim-5) and half-integer-spin
  (dim-2 spinor, dim-4, dim-6), matching the 2I irrep parity
  under the central element $-e$.
\item The full permutation representation $\chi_{\mathrm{perm}}$
  on 600 cells decomposes into 2I irreps with integer
  multiplicities $(5, 10, 10, 15, 15, 20, 20, 25, 30)$ on
  $(1, 2a, 2b, 3a, 3b, 4a, 4b, 5, 6)$, exactly $5 \cdot \dim \rho$
  per irrep — matching 5 copies of the regular $2I$ rep, as
  forced by the free action.
\end{itemize}

\paragraph{What is identified ($\triangleright$).}
\begin{itemize}
\item The $2I$-orbit decomposition of the cell graph has $5$
  orbits of size $120$ each (sim-verified by build B\_orbits at
  \texttt{scripts/wo4\_sim\_close\_orbits.py}, \emph{not} imported
  from WO-1; WO-1 catalogue R.0 supplies only the shared
  $2I/A_5$ substrate). Build B\_laplacian at
  \texttt{scripts/wo4\_sim\_close\_laplacian.py} further shows
  the $600 \times 600$ cell-graph Laplacian has
  low-lying eigenvalues of multiplicities $1, 4, 9, 16, 25, 36$
  — exactly the squared dimensions of the first six $2I$
  irreps, confirming isotypic structure on low modes.
\item The $A_5$ irrep count is $5$, giving at most $5$
  morphologically distinct attractor families under the conjecture.
\end{itemize}

\paragraph{What is open ($\circ$).}
\begin{itemize}
\item Construction of $\Lcasc$ (Definition~\ref{def:D2} is
  conjectural; three construction routes suggested).
\item Conjecture~\ref{conj:C1} (morphological attractors as
  eigenmodes): central framework claim.
\item Conjecture~\ref{conj:C2} (double-head / no-head as
  $\mathbf{3}$-inversion): specific match to planarian data.
\item Conjecture~\ref{conj:C3} (forbidden-pattern falsifier):
  the cleanest experimental test.
\item All empirical comparison against Levin-group data.
\end{itemize}

\paragraph{Status summary.} This is a framework conjecture
note rather than a substantive derivation paper. It exists in
programme-drop form to ensure the Levin-bioelectric prediction is
on the record alongside the other biology-rung papers, not to
claim a proved result. Its falsifiability is stated explicitly;
the computation and empirical comparison are flagged for future
work.

%==================================================================
\section{Reproducibility}\label{sec:repro}
%==================================================================

This manuscript lives at
\texttt{papers/biology-rung-bioelectric/biology-rung-bioelectric.tex};
umbrella-programme artefacts are in \texttt{papers/biology-rung/}
(WO spec at \texttt{WO-BIOLOGY-RUNG-004-BIOELECTRIC.md}; shared
math-catalogue ledger). No computational artefact is supplied by
this paper; the $\Lcasc$-eigenmode sim is flagged as
\texttt{WAITING\_FOR\_CONSTRUCTION} pending either resolution of
the Hopf cell-fibration or an alternative graph-theoretic
definition of $\Lcasc$.

\begin{thebibliography}{99}

\bibitem{LevinBeane2011}
Beane, W.~S., Morokuma, J., Adams, D.~S., and Levin, M.,
\emph{A chemical genetics approach reveals H,K-ATPase-mediated
membrane voltage is required for planarian head regeneration},
Chemistry \& Biology \textbf{18} (2011), 77--89.

\bibitem{DurantLevin2017}
Durant, F., Morokuma, J., Fields, C., Williams, K., Adams, D.~S.,
and Levin, M.,
\emph{Long-term, stochastic editing of regenerative anatomy via
targeting endogenous bioelectric gradients},
Biophysical Journal \textbf{112} (2017), 2231--2243.

\bibitem{EmmonsBell2015}
Emmons-Bell, M., Durant, F., Hammelman, J., Bessonov, N.,
Volpert, V., Morokuma, J., Pinet, K., Adams, D.~S., Pietak, A.,
Lobo, D., and Levin, M.,
\emph{Gap junctional blockade stochastically induces different
species-specific head anatomies in genetically wild-type
Girardia dorotocephala flatworms},
International Journal of Molecular Sciences \textbf{16} (2015),
27865--27896.

\bibitem{VandenbergLevin2011}
Vandenberg, L.~N., Morrie, R.~D., and Adams, D.~S.,
\emph{V-ATPase-dependent ectodermal voltage and pH regionalization
are required for craniofacial morphogenesis},
Developmental Dynamics \textbf{240} (2011), 1889--1904.

\bibitem{PietakLevin2016}
Pietak, A., and Levin, M.,
\emph{Exploring instructive physiological signaling with the
Bioelectric Tissue Simulation Engine},
Frontiers in Bioengineering and Biotechnology \textbf{4} (2016), 55.

\bibitem{Levin2021}
Levin, M.,
\emph{Bioelectric signaling: Reprogrammable circuits underlying
embryogenesis, regeneration, and cancer},
Cell \textbf{184} (2021), 1971--1989.

\bibitem{CascadeBio}
\emph{Cascade Biology Rung: The 600-Cell Cell-Level Substrate},
internal manuscript,
\texttt{papers/cascade-derivation/cascade-bio.md}. Preprint,
Institute of Vibrational Field Dynamics.

\bibitem{CascadeFoundations}
\emph{Cascade Foundations: F1--F8}, internal manuscript,
\texttt{papers/cascade-derivation/cascade-foundations.md}. Preprint,
Institute of Vibrational Field Dynamics.

\bibitem{SmartXXXII}
Smart, L.,
\emph{The Closure Functional from First Principles: Self-Similar
Permeability, the 600-Cell, and the Derived Interaction}
(Paper XXXII),
\texttt{papers/paper-xxxii/paper-xxxii.tex}. Preprint,
Institute of Vibrational Field Dynamics, 2026.

\bibitem{SmartCapsid}
Smart, L.,
\emph{Caspar--Klug Triangulation Numbers from a Face-Level
Cascade Operator} (WO-1),
\texttt{papers/biology-rung-capsid/biology-rung-capsid.tex}.
Preprint, Institute of Vibrational Field Dynamics, 2026.

\bibitem{SmartAmyloid}
Smart, L.,
\emph{Cross-$\beta$ Amyloid Fibril Twist from $2I$ Helical-Orbit
Classes} (WO-2),
\texttt{papers/biology-rung-amyloid/} (in preparation), 2026.

\bibitem{SmartTubulin}
Smart, L.,
\emph{The 13-Protofilament Microtubule as a $C_{13}$ Selector}
(WO-3),
\texttt{papers/biology-rung-tubulin/biology-rung-tubulin.tex}.
Preprint, Institute of Vibrational Field Dynamics, 2026.

\bibitem{SmartPhyllotaxis}
Smart, L.,
\emph{Fibonacci Phyllotaxis from the 600-Cell's $\phig$-Helical
Closure Orbits} (WO-5),
\texttt{papers/biology-rung-phyllotaxis/biology-rung-phyllotaxis.tex}.
Preprint, Institute of Vibrational Field Dynamics, 2026.

\end{thebibliography}

\end{document}
